\hypertarget{chapter-14-graphics-as-a-pipeline-not-an-attachment}{%
\chapter{Graphics as a Pipeline, Not an Attachment}\label{chapter-14-graphics-as-a-pipeline-not-an-attachment}}

The default way most people put a figure into a LaTeX document is \texttt{\textbackslash{}includegraphics}, pointing at an image file produced somewhere else, by some other program, and pasted in much the way a figure would be pasted into a word processor. This works, and for a photograph or a hand-drawn illustration it is the only sensible approach, since there is no meaningful sense in which such an image could be generated from the document's own source. For a plot, a diagram, or any figure that represents data or a formal structure the document itself already describes, treating the figure as an external attachment throws away one of LaTeX's genuine advantages: the ability to define the figure in the same source language as the surrounding text, so that regenerating the document also regenerates the figure, correctly, from whatever the current state of the underlying data or argument actually is.

\hypertarget{tikz-for-figures-defined-in-the-document}{%
\section{TikZ for Figures Defined in the Document}\label{tikz-for-figures-defined-in-the-document}}

TikZ is a graphics language embedded in LaTeX itself, letting a figure be specified as LaTeX source rather than as a separate binary image file. A diagram drawn in TikZ lives in the same file as the text that discusses it, uses the same fonts as the surrounding document automatically, since it is typeset by the same engine rather than imported as a pre-rendered image, and can reference the document's own macros and notation directly, so a diagram's labels stay consistent with the notation used in the prose around it without any separate effort to keep the two synchronized.

This matters most for diagrams that need to change as an argument develops. A hand-drawn or externally-produced diagram that needs updating requires opening a separate tool, making the edit, re-exporting, and re-importing; a TikZ diagram is edited exactly like any other part of the document's source, and its update is included in the same commit, the same diff, and the same build as everything else, which keeps a diagram from silently going stale relative to a revised argument the way an externally maintained image file easily can. For any figure that is fundamentally schematic, an architecture diagram, a state machine, a commutative diagram as covered in Chapter 11, a flowchart describing a process the surrounding text also describes, TikZ is the tool that keeps the figure and the argument it illustrates from drifting apart over the course of a long revision history.

\hypertarget{pgfplots-and-data-driven-figures}{%
\section{PGFPlots and Data-Driven Figures}\label{pgfplots-and-data-driven-figures}}

PGFPlots, built on top of TikZ, extends this same in-document approach to actual plots: line graphs, scatter plots, bar charts, and surface plots, specified either with inline coordinate data or, more usefully for a figure that represents real data rather than a hand-constructed illustration, by reading directly from an external data file at compile time. \texttt{\textbackslash{}addplot\ table\ \{data.csv\};} pulls plotting data straight from a CSV file, meaning the figure updates automatically the next time the document is built, if the underlying CSV file has changed, with no manual step of re-exporting a plot image and re-importing it.

This is the concrete mechanism behind the principle this chapter opens with: a figure built this way is not attached to the document, it is derived from the same data the document's own build process already has access to. A plot showing experimental results stays current with the latest data automatically; a plot showing results from an earlier stage of an analysis that has since been revised will visibly and correctly update the next time the document is rebuilt, rather than silently continuing to show stale results because nobody remembered to regenerate and re-import a separate image file.

\hypertarget{csv-to-table-workflows}{%
\section{CSV-to-Table Workflows}\label{csv-to-table-workflows}}

The same principle applies to tables as to plots, and is arguably even more valuable there, since a hand-transcribed table is one of the most common sources of a document silently diverging from its actual underlying data. \texttt{pgfplotstable}, part of the PGFPlots ecosystem, reads a CSV file and typesets it as a properly formatted LaTeX table directly, with support for column-specific number formatting, computed columns derived from the source data, and filtering or sorting applied at typesetting time rather than requiring the CSV itself to already be in exactly the right order and format.

For a book presenting any substantial amount of tabular data, results tables, comparison tables, parameter listings, generating these directly from the CSV or other structured data file that is the actual source of truth, rather than transcribing values into a hand-written \texttt{tabular} environment, removes an entire category of error: a hand-transcribed table can silently diverge from its source the moment the source is updated and the table is not, while a table generated at compile time from the same file cannot diverge from that file by construction, since it is reading it fresh on every build.

\hypertarget{computation-outside-latex-included-inside-it}{%
\section{Computation Outside LaTeX, Included Inside It}\label{computation-outside-latex-included-inside-it}}

Not every figure's data originates in a format PGFPlots or \texttt{pgfplotstable} can read directly, and not every plot is simple enough to be worth expressing in PGFPlots's own plotting language rather than a dedicated plotting library in Python, R, or Julia. For these cases, the pipeline principle still applies, just with an external step made explicit rather than avoided: a script, checked into the project alongside the document's own source as described in Chapter 4, reads the underlying data, produces a plot as a PDF or SVG, and the document includes that output with \texttt{\textbackslash{}includegraphics} exactly as it would any other image, with the crucial difference that the generating script is part of the project, versioned alongside everything else, and rerunning it is a documented, repeatable step rather than a one-off action performed once and then forgotten.

Wiring this generation step into the build process itself, through a Makefile target as described in Chapter 5 that runs the plotting script before invoking \texttt{latexmk}, or by making the figure a dependency the Makefile knows how to regenerate whenever its source data changes, closes the remaining gap between this hybrid workflow and the fully in-document approach of TikZ and PGFPlots: even when the actual plotting happens outside LaTeX, ``rebuild the document'' can still mean ``regenerate every figure from current data and then typeset,'' rather than ``typeset whatever figure files happen to already be sitting in the figures directory,'' which is precisely the guarantee that makes a figure trustworthy months or years after it was first produced.

The next chapter extends this same pipeline further, bringing computation more directly inside the document itself: embedded code listings, verified output, and simulation results woven into the text as a single build rather than a document written up separately from the computation it describes.
